\chapter{Technical notes}\label{ch:technical}

This chapter is for users who fit large models, who run \nwsaom{} on several platforms, or who want to read or extend the source. It describes the three layers of the program, the native plugin and how to tell which engine ran, the threads, the size limits of Stata and of \nwsaom, the run times to expect, and the limit on the size of an ado program that shapes the layout of the source files.

\section{The three layers}

\nwsaom{} consists of three layers.
\begin{enumerate}[nosep]
  \item The \textbf{ado layer} (\texttt{nwsaom.ado} and a few helper files) parses the options, checks the data and the combinations of options, builds the model description, calls the estimator and posts the results in \texttt{e()}.
  \item The \textbf{Mata core} (\texttt{unw\_saom.do}, compiled into the library \texttt{lnwcommands.mlib}) contains the model description, the observed statistics, the Robbins--Monro estimator for all phases, the calculation of the covariance matrix and a reference simulator, written in Mata, for the core effects.
  \item {\sloppy The \textbf{native plugin} (\texttt{native/saom\_sim.c}, compiled for each platform into \texttt{lib/plugins/<platform>/}\allowbreak\texttt{saom\_sim\_<platform>.plugin}) simulates the ministeps. Almost all of the computing time of a fit is spent in it.\par}
\end{enumerate}
The Mata layer talks to the plugin through Stata frames. The calls are made once per batch of simulations and not once per ministep: in phases 1 and 3 all replicates of a fit are run in one call, on worker threads. The ministep sampler itself is the multinomial-logit rule of Chapter~\ref{ch:model}: for the acting actor the contribution of every alternative to the objective function is computed, the alternatives are given the choice probabilities, and one is drawn.

\section{The native plugin}

\subsection{Protocol versions}

The plugin and the Mata layer communicate through a wire protocol. The plugin reports the version of the protocol that it implements, and the Mata layer checks the version before it uses a feature that needs a newer one. The version of the plugin in the repository at the time of writing is 62. New effects, and new ways of combining existing ones, increase the version; for example, version 11 introduced network and behavior endowment/creation, structural values and the threaded path for non-directed models and rate covariates, versions 26 and 27 composition change at fractional times, version 52 several co-evolving networks, and version 56 covariate effects in two-mode networks.

A plugin that is older than the \texttt{.ado} files cannot run models that need a newer protocol. This happens most often on Linux and Windows: those two binaries are built by continuous integration (a workflow on GitHub), whereas the macOS binary is compiled by the developer, so that the three can differ for a while after a commit. If you install from a repository checkout, \nwsaom{} looks for the plugin first in \texttt{lib/plugins/<platform>/} next to the \texttt{nwsaom.ado} that is being run and only then with \texttt{findfile}. After a \texttt{net install} there is no \texttt{lib/plugins/} next to \texttt{nwsaom.ado}, and the copy in the PLUS directory is used.

\subsection{Native-only effects}

Not all effects are implemented twice. The Mata layer contains a simulator for the older core of the model, which is used if no plugin is available for the platform. The effects added later, for the most part the large families of RSiena effects that were ported after the core was finished, exist only in the plugin: the plugin computes the simulation and the observed statistic, and there is no Mata twin. This is a design decision and not an omission to be repaired; extending the C plugin is how new effects are added. In the help file such effects are marked \emph{native plugin only}, together with the protocol version in which they were introduced (for example, the covariate-by-triad effects of \opt{samextranstrip()} and its relatives, protocol 39).

If you ask for a native-only effect and the native engine cannot be used, \nwsaom{} stops with an error and does not fall back. With the Mata engine forced by the debugging switch described below, the request
\begin{lstlisting}
. nwsaom, wave1(glasgow1) wave2(glasgow2) outdegree reciprocity egoplusaltsqx(alcohol1) seed(1)
\end{lstlisting}
gives
\begin{lstlisting}
nwsaom: egoplusaltsqx needs the native plugin (protocol 41 or later); there is no Mata implementation of this effect.
\end{lstlisting}
and the return code is 198. The reverse also holds: a model that consists of effects with a Mata implementation runs without a plugin, correctly but slowly.

Two consequences follow for the user. First, with a plugin that is too old the same model is either slower (Mata) or refused (native-only), and the message says which. Second, to use a recent effect on Linux or Windows you may have to wait for, or build, a current plugin; the C source and the build rules are in the repository.

\subsection{Which engine ran}

After a fit, \texttt{e(engine)} is \texttt{native} or \texttt{mata}. If it is \texttt{mata}, a note is printed after the table and \texttt{e(engine\_why)} gives the reason. Look at it after every large fit. The Mata engine gives statistically equivalent results (the random streams of the two engines are not the same, so that the digits differ), but it is slower by orders of magnitude. As an illustration, a two-wave Glasgow model with outdegree and reciprocity and \texttt{k0(10) k3(10)} took about 30 seconds in the Mata engine, where the native engine needs a small fraction of a second for the same model.

For tests, setting the global macro \texttt{NWSAOM\_NONATIVE} to 1 before the call forces the Mata engine (and \texttt{e(engine\_why)} then gives the message of a missing plugin); setting \texttt{NWSAOM\_TRACE} to 4 prints a line with the protocol version of the plugin that was found and whether the native and the batch paths are used. These are development switches.

\subsection{Fixed limits of the plugin}

The plugin has fixed capacities: at most 256 effects in a model, 1\,024 covariate arrays (a dyadic covariate on $n$ actors uses $n$ of them, and its usable limit is quoted as 1\,023), 64 behavior effects, 8 co-evolving behaviors and 16 rate covariates. A model that exceeds one of them runs in the Mata engine without an error if the effects allow it, which is a trap discussed in Section~\ref{sec:limits}.

\section{Threads and reproducibility}\label{sec:threads}

The plugin is multi-threaded. In phase 1 and phase 3 all replicates of all periods are run in one call by a pool of worker threads. Phase 2 is sequential by construction, because each Robbins--Monro step needs the previous one; there only the periods of one step are run in parallel, and only when the network is large enough for threads to pay off.

\paragraph{The number of threads.} \opt{cores(\#)} sets it; the default 0 uses all physical cores and \opt{cores(1)} runs single-threaded.

\paragraph{Reproducibility.} Every simulated period draws from its own random stream, which is derived from the seed and from its position (replicate and period), and sums over the periods are formed after the threads have joined, in a fixed order. A given \opt{seed()} therefore gives identical results whatever the number of threads. This was checked in the preparation of this manual on the two-wave example with \opt{cores(1)} and \opt{cores(4)}: the coefficients were identical to all printed digits (the outdegree coefficient in both runs $-2.2380415$). The results change if the model, the seed or the version of the program changes.

\section{Size limits}\label{sec:limits}

\subsection{What Stata limits}

Everything that is handed to \nwsaom{} as a Stata matrix is limited by the edition of Stata:

\begin{center}
\small
\begin{tabular}{@{}llll@{}}
\toprule
 & Stata/BE & Stata/SE & Stata/MP \\
\midrule
Largest matrix & $800\times 800$ & $11\,000\times 11\,000$ & $65\,534\times 65\,534$ \\
Variables & 2\,048 & 32\,767 & 120\,000 \\
\bottomrule
\end{tabular}
\end{center}
Only the Stata/BE values were checked on the development machine; the SE and MP values are taken from the Stata documentation. The consequences for networks are these. A network that is declared from a matrix (\texttt{nwset, mat()}), and every matrix given to \opt{dyadcov()}, \opt{dyadcovrecip()}, \opt{dyadx()}, \opt{missnet()} or \opt{structural()}, is limited by the matrix size: 800 actors in Stata/BE. Larger networks are brought in as edge lists (\texttt{nwfromedge}); nothing dense is needed for the network itself. A dyadic covariate also travels to the plugin as one variable per actor, so that the native engine handles it up to about 1\,000 actors (1\,023 covariate arrays), and Stata/BE with its 2\,048 variables stops at about 2\,000 actors in any case. Mata itself has no such limits; the networks are kept as sparse adjacency lists, and only dense by-products scale with $n^2$.

\subsection{What \nwsaom{} costs}

Every simulated ministep considers all $n$ possible alters, and the number of ministeps per period grows with $n$ (the rate times $n$), so that the time of a fit grows roughly with $n^3$ above about 1\,000 actors. For a plain model (outdegree, reciprocity, transitive triplets; three waves; about ten ties per actor; native engine; \texttt{k0(30) k3(200)}) the time was measured on one machine as follows.

\begin{center}
\small
\begin{tabular}{@{}rr@{}}
\toprule
Actors & Seconds \\
\midrule
250 & 1.5 \\
500 & 2.6 \\
1\,000 & 10 \\
2\,000 & 67 \\
4\,000 & 520 \\
\bottomrule
\end{tabular}
\end{center}
A full-length run, with the default \opt{k3()} of 1\,000 and more phase-2 iterations, costs several times more. One thousand actors is routine; a few thousand actors take minutes to tens of minutes; beyond roughly 10\,000 actors the model itself, in which each ministep looks at all alters, is the limit and not the software. A dyadic covariate adds to this (at $n=300$ about 4 seconds instead of 2, at $n=700$ about 12 seconds in the measurement quoted).

\subsection{The silent fallback}

\textbf{If a model exceeds a capacity of the native engine but all its effects have a Mata implementation, the fit runs in Mata without an error.} It gives the same results, but it is slower by orders of magnitude. Until 22 October 2026 the test for ``too many covariate arrays'' used 127 instead of the plugin's real limit of 1\,023, so that a dyadic covariate on 129 or more actors ran in Mata: a fit that takes seconds natively ran for more than an hour. That is fixed, but the lesson holds: after you have fitted a large network, look at \texttt{e(engine)}.

\section{Performance}

The following timings were measured on one 18-core machine with the native engine, and are given as orders of magnitude (the figures of the development notes).

\begin{center}
\small
\begin{tabular}{@{}lr@{}}
\toprule
Model & Seconds \\
\midrule
Glasgow, 3 waves, outdegree, reciprocity, transitive triplets & 0.33 \\
Glasgow, 2 waves, outdegree, reciprocity & 0.23 \\
Glasgow, 3 waves, co-evolution with alcohol & 0.63 \\
Synthetic, 500 actors, 3 waves, co-evolution & 18 \\
\bottomrule
\end{tabular}
\end{center}
For comparison, the larger Glasgow model of the validation (with \texttt{gwespFF}, \texttt{transRecTrip}, \texttt{inPopSqrt}, \texttt{outAct}, covariate effects and a rate covariate) took \nwsaom{} about 1 second against 5--7 seconds for \texttt{siena07}. The examples in Chapter~\ref{ch:start} took one to a few seconds each on the development machine, including the start of Stata.

Some hints. Use the default number of cores. Shorten exploratory runs with \opt{k0()} and \opt{k3()} (as the example fragments of Chapter~\ref{ch:options} do) and use the defaults for the final run, because the standard errors and the convergence ratios come from phase 3. Models with many weakly identified effects take more iterations, not less.

\section{The size limit of an ado program}

Stata refuses to load a single ado program larger than about 141\,000 bytes of code (comments and blank lines are not counted). The limit is a property of Stata and not of the file size: a program that is only a few lines above it fails to load with the message ``system limit exceeded -- see manual (error occurred while loading nwsaom.ado)'', also when the same file loads in an empty session. The main program \texttt{nwsaom} is close to this cap, and the layout of the source is shaped by it.

If you modify the ado code, remember the rules that were learnt:
\begin{itemize}[nosep]
  \item Move blocks of code into separate programs in \emph{separate ado-files}. The programs of one ado-file all count towards Stata's total adosize of 1\,000\,KB while loaded, and a program that another ado-file calls must have its own file (a program defined inside \texttt{nwsaom.ado} can be evicted when the adosize fills up).
  \item Mata functions defined in an ado-file are private to that file; helpers that several ado-files share belong in \texttt{unw\_saom.do}.
  \item A subprogram receives the locals of its caller only through options. Results come back in \texttt{r()}.
  \item Do not name a new option as an extension of an existing one (\texttt{netx0} next to \texttt{netx}): Stata's \texttt{syntax} routes the shorter name to the option declared first. Option names that are a digit followed by letters are rejected altogether, which is why the options \opt{dblnbrdist2} and \opt{inplus3} and the arguments of \texttt{nwsaom multiplex} look as they do.
\end{itemize}

\section{File layout}

\begin{center}
\small
\begin{tabular}{@{}>{\raggedright\arraybackslash}p{5.2cm}>{\raggedright\arraybackslash}p{9.8cm}@{}}
\toprule
File & Content \\
\midrule
\texttt{nwsaom.ado} & The command: parsing, checks, model description, estimation dispatch, results. \\
\texttt{\_nwsaom\_nx.ado} & The native-only network options \opt{netx()}, \opt{ownx()}, \opt{dyadx()}, \opt{tvdyadcov()}, \opt{tvdyadcovrecip()}, \opt{tvdyadx()}, moved out of \texttt{nwsaom.ado} because of the size limit. \\
\texttt{nwsaom\_multinet.ado} & \texttt{nwsaom multinet}: several dependent networks. \\
\texttt{nwsaom\_multiplex.ado} & \texttt{nwsaom multiplex}: the earlier two-network subcommand with outdegree, reciprocity and the cross-network product. \\
\texttt{\_nwsaom\_postengine.ado}, \texttt{\_nwsaom\_rsnotes.ado}, \texttt{\_nwsaom\_tconvtable.ado} & Posting of \texttt{e(engine)}, the notes with RSiena's effect names, and the printing of the convergence table. \\
\texttt{nwsaom\_estat.ado}, \texttt{nwsaom\_meta.ado} & \texttt{estat gof} and \texttt{estat mems}; the meta-analysis. \\
\texttt{unw\_saom.do} & The Mata core: model, estimator, reference simulator, wire protocol. \\
\texttt{native/saom\_sim.c} & The C plugin source; \texttt{native/Makefile} builds it. \\
\texttt{lib/plugins/<platform>/} & The compiled plugins for \texttt{macos}, \texttt{unix} and \texttt{windows}. \\
\texttt{nwsaom.sthlp}, \texttt{nwsaom\_remarks.sthlp} & The help file and its remarks. The help file was split in two because Stata's help Viewer fails on a very long file. \\
\texttt{cscripts/test\_nwsaom\_*.do} & The regression tests, one file per feature, and \texttt{cscripts/run\_nwsaom\_tests.sh}, which runs them (\texttt{quick} or \texttt{full}) and accepts the directory of a freshly built plugin. \\
\bottomrule
\end{tabular}
\end{center}

\paragraph{Testing.} The tests are run in two modes: in \emph{development mode}, in which the Mata sources are loaded directly (\texttt{do unw\_core.do}, \texttt{do unw\_ergm.do}, \texttt{do unw\_saom.do}, and \texttt{nwclear} --- never \texttt{clear all}, which would delete the loaded Mata functions) and in \emph{production mode} against the compiled library and plugin. The native simulator is checked against the Mata simulator at fixed parameters for statistical equivalence of the simulated statistics (the two do not share a random stream), and the estimates are checked against RSiena as described in the validation chapter.

\paragraph{A stale plugin.} A plugin from an earlier installation in the PLUS directory can shadow the plugin of a freshly built checkout in some set-ups and silently run an old binary. If a new effect is refused as unknown to the plugin, or the protocol shown by \texttt{NWSAOM\_TRACE} is older than expected, check which plugin file is used.
